% GENERATED FILE. Edit canonical lessons, then run npm run generate:books.
% canonical-source-hash: fnv1a64:7f9baddcd491dcad
% canonical-lessons: MR-C553-sadhya-karne, MR-C553-hulkavni-dene, MR-C553-acuk, MR-C553-hamkhas, MR-C553-tatpar

\chapter{Reaching the Goal}
\label{ch:mr-reaching-the-goal}
\hlchaptermodality{553}

\begin{chapteropening}
\textbf{By the end of this chapter:} \emph{I can talk about achieving aims, accuracy and being ready.}

Complete the last lesson of chapter 553: I can talk about achieving aims, accuracy and being ready.
\end{chapteropening}

\section[\texorpdfstring{\mr{साध्य} \mr{करणे} (sādhya karṇe)}{साध्य करणे (sādhya karṇe)}]{\mr{साध्य} \mr{करणे} (sādhya karṇe) — to achieve}

\label{lesson:MR-C553-sadhya-karne}



\begin{quote}
Before the new one: say the Marathi for to break down, then the Marathi for to come loose.
\end{quote}

\subsection*{You'll want to know: \mr{साध्य} \mr{करणे}}
\textbf{\mr{साध्य} \mr{करणे}} — \emph{sādhya karṇe} — \textquotedblleft{}to achieve\textquotedblright{}.

\begin{quote}
\mr{खूप} \mr{कष्ट} \mr{करून} \mr{तिने} \mr{तिचं} \mr{ध्येय} \mr{साध्य} \mr{केलं}.
\end{quote}

\begin{quote}
\emph{khūp kaṣṭa karūn tine ticaṁ dhyey sādhya kelaṁ.}
\end{quote}

\begin{quote}
With much hard work she achieved her goal.
\end{quote}

Past, with -\mr{ने} (he …): \mr{त्याने} \mr{ध्येय} \mr{साध्य} \mr{केलं}, \emph{tyāne dhyey sādhya kelaṁ}.

\begin{cousinweb}
\textbf{\mr{साध्य} \mr{करणे}} (\emph{sādhya karṇe}) is the tatsama \mr{साध्य}, from Sanskrit \emph{sādhya}, \textquotedblleft{}to be accomplished\textquotedblright{} (Monier-Williams), with \mr{करणे}.
\end{cousinweb}

\begin{grammarlens}[title={What you've built}]
One more word for saying what happens if.
\end{grammarlens}

\subsection*{Your turn}
\begin{itemize}
\raggedright
  \item \emph{Say it:} \emph{sādhya karṇe}
  \item \emph{Say it:} \emph{sādhya karṇe}, once more
  \item \emph{Say it:} \emph{khūp kaṣṭa karūn tine ticaṁ dhyey sādhya kelaṁ.} — With much hard work she achieved her goal.
  \item \emph{Recall:} say the Marathi for to break down, then the Marathi for to come loose, then say \emph{sādhya karṇe} again
\end{itemize}

\begin{tcolorbox}[breakable,colback=teal!4,colframe=teal!35!black,title={Before you move on}]
What does \mr{साध्य} \mr{करणे} mean? (\textquotedblleft{}To achieve\textquotedblright{}.) Say it once more.
\end{tcolorbox}
\section[\texorpdfstring{\mr{हुलकावणी} \mr{देणे} (hulkāvṇī deṇe)}{हुलकावणी देणे (hulkāvṇī deṇe)}]{\mr{हुलकावणी} \mr{देणे} (hulkāvṇī deṇe) — to elude; to give the slip}

\label{lesson:MR-C553-hulkavni-dene}



\begin{quote}
Before the new one: say the Marathi for to come loose, then the Marathi for to achieve.
\end{quote}

\subsection*{You'll want to know: \mr{हुलकावणी} \mr{देणे}}
\textbf{\mr{हुलकावणी} \mr{देणे}} — \emph{hulkāvṇī deṇe} — \textquotedblleft{}to elude; to give the slip\textquotedblright{}.

\begin{quote}
\mr{यशाने} \mr{त्याला} \mr{दोनदा} \mr{हुलकावणी} \mr{दिली}.
\end{quote}

\begin{quote}
\emph{yaśāne tyālā dondā hulkāvṇī dilī.}
\end{quote}

\begin{quote}
Success eluded him twice.
\end{quote}

Past, with -\mr{ने} (he …): \mr{त्याने} \mr{हुलकावणी} \mr{दिली}, \emph{tyāne hulkāvṇī dilī}.

The one dodged takes -\mr{ला}: \mr{त्याला} \mr{हुलकावणी} \mr{दिली}.

\begin{cousinweb}
\textbf{\mr{हुलकावणी} \mr{देणे}} (\emph{hulkāvṇī deṇe}) is \mr{हुलकावणी}, \textquotedblleft{}a dodge, a feint\textquotedblright{}, with \mr{देणे}; the origin of \mr{हुलकावणी} is not settled.
\end{cousinweb}

\begin{grammarlens}[title={What you've built}]
One more word for saying what happens if.
\end{grammarlens}

\subsection*{Your turn}
\begin{itemize}
\raggedright
  \item \emph{Say it:} \emph{hulkāvṇī deṇe}
  \item \emph{Say it:} \emph{hulkāvṇī deṇe}, once more
  \item \emph{Say it:} \emph{yaśāne tyālā dondā hulkāvṇī dilī.} — Success eluded him twice.
  \item \emph{Recall:} say the Marathi for to come loose, then the Marathi for to achieve, then say \emph{hulkāvṇī deṇe} again
\end{itemize}

\begin{tcolorbox}[breakable,colback=teal!4,colframe=teal!35!black,title={Before you move on}]
What does \mr{हुलकावणी} \mr{देणे} mean? (\textquotedblleft{}To elude; to give the slip\textquotedblright{}.) Say it once more.
\end{tcolorbox}
\section[\texorpdfstring{\mr{अचूक} (acūk)}{अचूक (acūk)}]{\mr{अचूक} (acūk) — accurate; unerring}

\label{lesson:MR-C553-acuk}



\begin{quote}
Before the new one: say the Marathi for to achieve, then the Marathi for to elude.
\end{quote}

\subsection*{You'll want to know: \mr{अचूक}}
\textbf{\mr{अचूक}} — \emph{acūk} — \textquotedblleft{}accurate; unerring\textquotedblright{}.

\begin{quote}
\mr{त्याचं} \mr{उत्तर} \mr{अगदी} \mr{अचूक} \mr{होतं}.
\end{quote}

\begin{quote}
\emph{tyācaṁ uttar agdī acūk hotaṁ.}
\end{quote}

\begin{quote}
His answer was quite accurate.
\end{quote}

Built on \mr{चूक}, a mistake: without a mistake.

\begin{cousinweb}
\textbf{\mr{अचूक}} (\emph{acūk}) is built on \mr{चूक}, \textquotedblleft{}a mistake\textquotedblright{}, with \emph{a-}, \textquotedblleft{}without\textquotedblright{}; CDIAL lists Marathi \emph{cukṇẽ}, \textquotedblleft{}to blunder\textquotedblright{}, under \emph{*cukk} (CDIAL 4848).
\end{cousinweb}

\begin{grammarlens}[title={What you've built}]
One more word for saying what happens if.
\end{grammarlens}

\subsection*{Your turn}
\begin{itemize}
\raggedright
  \item \emph{Say it:} \emph{acūk}
  \item \emph{Say it:} \emph{acūk}, once more
  \item \emph{Say it:} \emph{tyācaṁ uttar agdī acūk hotaṁ.} — His answer was quite accurate.
  \item \emph{Recall:} say the Marathi for to achieve, then the Marathi for to elude, then say \emph{acūk} again
\end{itemize}

\begin{tcolorbox}[breakable,colback=teal!4,colframe=teal!35!black,title={Before you move on}]
What does \mr{अचूक} mean? (\textquotedblleft{}Accurate; unerring\textquotedblright{}.) Say it once more.
\end{tcolorbox}
\section[\texorpdfstring{\mr{हमखास} (hamkhās)}{हमखास (hamkhās)}]{\mr{हमखास} (hamkhās) — without fail; surely}

\label{lesson:MR-C553-hamkhas}



\begin{quote}
Before the new one: say the Marathi for to elude, then the Marathi for accurate.
\end{quote}

\subsection*{You'll want to know: \mr{हमखास}}
\textbf{\mr{हमखास}} — \emph{hamkhās} — \textquotedblleft{}without fail; surely\textquotedblright{}.

\begin{quote}
\mr{तो} \mr{दर} \mr{रविवारी} \mr{हमखास} \mr{आईला} \mr{फोन} \mr{करतो}.
\end{quote}

\begin{quote}
\emph{to dar ravivārī hamkhās āīlā phon karto.}
\end{quote}

\begin{quote}
He calls his mother every Sunday without fail.
\end{quote}

\begin{cousinweb}
\textbf{\mr{हमखास}} (\emph{hamkhās}) is a loan from Persian, built from \emph{ham}, \textquotedblleft{}also\textquotedblright{}, and Arabic \emph{khāṣṣ}, \textquotedblleft{}special\textquotedblright{}.
\end{cousinweb}

\begin{grammarlens}[title={What you've built}]
One more word for saying what happens if.
\end{grammarlens}

\subsection*{Your turn}
\begin{itemize}
\raggedright
  \item \emph{Say it:} \emph{hamkhās}
  \item \emph{Say it:} \emph{hamkhās}, once more
  \item \emph{Say it:} \emph{to dar ravivārī hamkhās āīlā phon karto.} — He calls his mother every Sunday without fail.
  \item \emph{Recall:} say the Marathi for to elude, then the Marathi for accurate, then say \emph{hamkhās} again
\end{itemize}

\begin{tcolorbox}[breakable,colback=teal!4,colframe=teal!35!black,title={Before you move on}]
What does \mr{हमखास} mean? (\textquotedblleft{}Without fail; surely\textquotedblright{}.) Say it once more.
\end{tcolorbox}
\section[\texorpdfstring{\mr{तत्पर} (tatpar)}{तत्पर (tatpar)}]{\mr{तत्पर} (tatpar) — prompt; ready to act}

\label{lesson:MR-C553-tatpar}



\begin{quote}
Before the new one: say the Marathi for accurate, then the Marathi for without fail.
\end{quote}

\subsection*{You'll want to know: \mr{तत्पर}}
\textbf{\mr{तत्पर}} — \emph{tatpar} — \textquotedblleft{}prompt; ready to act\textquotedblright{}.

\begin{quote}
\mr{हा} \mr{दुकानदार} \mr{ग्राहकांच्या} \mr{सेवेसाठी} \mr{नेहमी} \mr{तत्पर} \mr{असतो}.
\end{quote}

\begin{quote}
\emph{hā dukāndār grāhkāñcyā sevesāṭhī nehmī tatpar asto.}
\end{quote}

\begin{quote}
This shopkeeper is always ready to serve his customers.
\end{quote}

\begin{cousinweb}
\textbf{\mr{तत्पर}} (\emph{tatpar}) is a tatsama, taken from Sanskrit \emph{tatpara}, \textquotedblleft{}intent on that\textquotedblright{} (Monier-Williams).
\end{cousinweb}

\begin{grammarlens}[title={What you've built}]
One more word for saying what happens if.
\end{grammarlens}

\subsection*{Your turn}
\begin{itemize}
\raggedright
  \item \emph{Say it:} \emph{tatpar}
  \item \emph{Say it:} \emph{tatpar}, once more
  \item \emph{Say it:} \emph{hā dukāndār grāhkāñcyā sevesāṭhī nehmī tatpar asto.} — This shopkeeper is always ready to serve his customers.
  \item \emph{Recall:} say the Marathi for accurate, then the Marathi for without fail, then say \emph{tatpar} again
\end{itemize}

\begin{tcolorbox}[breakable,colback=teal!4,colframe=teal!35!black,title={Before you move on}]
What does \mr{तत्पर} mean? (\textquotedblleft{}Prompt; ready to act\textquotedblright{}.) Say it once more.
\end{tcolorbox}
